\section{A Morality of Asceticism and Action}

In a continual effort to show what well-bred men used to think was normal and healthy, I will share the thoughts of a 19\textsuperscript{th} century Frenchman who was still attached to that tradition. Unlike those who think they discovered sliced bread for the first time a few weeks ago, we accept that there is nothing new under the sun. Although \textbf{Rene Guenon} offers an interesting way of presenting and relating certain ideas, there is nothing new. Anyone who believes that following a higher path ``atrophies'' any part of himself has wildly misunderstood the writings; rather it makes life better, your mind is freed from inner turmoil, your work becomes creative play.

Although action follows thought, practice precedes theory. Without spiritual practice and the development of the proper inner attitude, the intellectual part of Tradition cannot be understood. Actually, practice is mandatory and theory is optional, although most men believe the opposite.

\textbf{Abbé Henri de Tourville}\footnote{\url{https://en.wikipedia.org/wiki/Henri\_de\_Tourville}} was born in Paris in 1842 and died at the Castle of Tourville in 1903. Some of his letters were collected in the small volume under the title \textit{Letters of Direction}. What follows are some extracts on topics that have come up recently.

\paragraph{On Present Times and Complaining about them}


\begin{quotationx}
These present times are really not at all bad though some old-fashioned ideas are apt to make us see the progress of the world as if everything were upside down.

Most good and even saintly people have the bad habit of incessantly lamenting inwardly over everything under the sun. I do not know of what Paradise on earth they can have dreamt, the remembrance of which makes them despair, for such a Paradise certainly does not exist.
\end{quotationx}


\paragraph{Germinating in Secret}


\begin{quotex}
Here in this world we make many generous and intelligent attempts, yet to all appearance they are usually unsuccessful. At any rate their successes are small in proportion to their cost. Yet in reality the results for which they pave the way are enormous. Look for instance at the Roman Empire becoming entirely Christian on the very morrow of the most cruel persecutions; and this because everywhere \emph{the seeds had been germinating in secret.}
\end{quotex}


\paragraph{Qualifying for the Struggle}


\begin{quotationx}
It is not right to groan over the state of the world as if it were lost. What is actually happening is a clash between the old spirit and the new, a clash which is especially noticeable because the old spirit is realing how old it is and how nothing is looked at any longer from its point of view.

Most people are like sheep and follow, without much satisfaction to themselves, the lines of past tradition. A very small minority emerges, with great hesitation and amidst endless discussion to be faced by troublesome and pressing contradictions. \emph{It is however of the minority that you must be, when God has put you there by interior vocation and natural aptitude}.
\end{quotationx}


\paragraph{The World is Speeding Up}


\begin{quotex}
All human institutions will in the future show less and less vitality and flexibility. In any case they generally develop on quite different lines form those originally intended. History, when we come to study it, is full of such instances. But it is more especially true of the new and indeterminate world of the future, in which the speeding up of everything makes human calculations over long periods impossible.
\end{quotex}


\paragraph{The Forerunners}


\begin{quotex}
In every age God has scattered forerunners in the world. They are those who are ahead of their time and whose personal action is based on an inward knowledge of that which is to come. If you and I should happen to be forerunners, let us bless God for it, even though, living a century or two too soon, we may feel ourselves to be strangers in a foreign land.
\end{quotex}


\paragraph{Spiritual Independence}


\begin{quotex}
Be content with training yourself in spiritual independence and let all your personal conduct be in keeping with such independence. This is the really essential point. Do not be surprised if you do not win others over to your way of thinking. Those in whom any new truth is born have always to carry it for a long time in solitude.

Rejoice then in the light which you have been given and do not be surprised that it is so difficult to pass it on to others. It really is making its way, not so much through you or me as through force of circumstance. You are simply ahead of your time; it is a good thing to have long sight and to let your soul be illumined as soon as you are aware of the light.
\end{quotex}


\paragraph{On Being Yourself}


\begin{quotex}
God is working on you in order that you may no longer be a child, tossed about by every wind, a prey to external influences. He has given you your own grace, you won nature (insofar as it is good), your own distinctive character. You are therefore required to be yourself and not anyone else.

Live according to your own nature; inwardly without restriction; outwardly insofar as external conditions permit. I would compare you to a sailor he has not doubts as to the port for which he is making and if he is obliged to tack, it is in order to make his port. It is not that he has changed his mind; on the contrary, he makes of the changing winds in order that their very changeableness may bring him to his port.
\end{quotex}


\paragraph{A Divine Power Has Descended}


On a different tack, and as a forerunner, there is \textbf{Seneca}'s 41st letter to Lucillius, written during the decadent time of Nero. He describes a man capable of transcending the times he lives in, a period of brutalization and licentiousness.

\begin{quotationx}
We push one another into vice. And how can a man be recalled to salvation, when he has no one to retrain him and all mankind to urge him on?

If you see a man who is unterrified in the midst of dangers, untouched by desires, happy in adversity, peaceful amid the storm, who looks down upon men from a higher plane and view the gods on a footing of equality, will not a feeling of reverence for him steal over you? Will you not say:

``This quality is too great and too lofty to be regarded as resembling this petty body in which it dwells. \emph{A divine power has descended upon that man}.''

When a soul rises superior to other souls, when it is under control, when it passes through every experience as if it were of small account, when it smiles at our fears and at our prayers, it is stirred by a force from heaven. A thing like this cannot stand upright unless it be propped by the divine.

Therefore, a greater part of it abides in that place from whence it came down to earth. Just as the rays of the sun do indeed touch the earth but still abide at the source from which they are sent, even so the great and hallowed soul, which has come down in order that we may have a nearer knowledge of divinity, does indeed associate with us, but still cleaves to its origin; on that source it depends, thither it turns its gave and strives to go, and it concerns itself with our doings only as a being superior to ourselves.
\end{quotationx}


\paragraph{Ascetic Morality and a Morality of Action}


\textbf{Carl Jung} regards such an attitude as the harbinger of things to come, a leavening agent, as it were. The ancients held to the idea of a mediator in whose name new ways of love would be opened. That idea became a fact in Christianity, so that human society took an immense stride forward.

\begin{quotex}
This was not the result of any speculative, sophisticated philosophy, both of an elementary need in the great masses of humanity vegetating in spiritual darkness. They were evidently driven to it by the profoundest inner necessities, for humanity does not thrive in a state of licentiousness.
\end{quotex}


Interestingly, Jung claims that \textbf{Christianity} and \textbf{Mithraism} have the same meaning: the moral subjugation of the animal instincts. Mithras was the Logos emanated by God. Thus those who see some wide gulf between those two spiritual movements probably understand neither. That is because men today no longer understand the inner reasons for their rise. Perhaps these final words from Jung will awaken something within to recreate a tradition to follow.

\begin{quotex}
Both religions i.e., Christianity and Mithraism teach a distinctly ascetic morality and a morality of action. The latter is particularly true of Mithraism. Cumont says that Mithraism owed its success to the value of its morality which above all things favoured action. The followers of Mithras formed a sacred army in the fight against evil and among them were \emph{virgines} (nuns) and \emph{continentes} (ascetics).
\end{quotex}



\flrightit{Posted on 2013-10-15 by Cologero}

\begin{center}* * *\end{center}

\begin{footnotesize}
\begin{sffamily}
\texttt{X on 2013-10-16 at 00:44 said:}

``What is actually happening is a clash between the old spirit and the new, a clash which is especially noticeable because the old spirit is realing how old it is and how nothing is looked at any longer from its point of view.''

The true newest spirit is older than the oldest. ``Before Abraham Was, I Am.''

\hfill

\texttt{Thorsten on 2013-10-16 at 11:13 said:}

Excellent advice. De Tourville's style hearkens back to Marcus Aurelius but is markedly more upbeat. Mithraism is certainly inspirational in many regards and Evola's pamphlet on the subject was eye-opening if a bit redundant. To me it seems increasingly clear that a new Western religion would do well to integrate Platonic metaphysics as has been suggested, with the use of Indo-European beast-slaying motifs as seen in Mithraism, the Volsunga Saga, medieval legends of Saint Michael and St. George. etc. If society does undergo some drastic simplifications in the near future such a cult of heroes would be very apt methinks.

\hfill

\texttt{Jacob on 2013-10-16 at 13:33 said:}

Great post.

I dropped Jung after reading that he was untraditional, but I'm going to try reading him with a different eye after some of your recent comments.

\hfill

\texttt{scardanelli on 2013-10-16 at 14:33 said:}

@ Thorsten, I think that maybe it's not about the beast slaying motif itself, but the truth behind the myth- ``the moral subjugation of the animal instincts.'' If we couple this truth with Platonic metaphysics, then this ``religion of the future'' begins to sound strangely like Christianity.

``Unlike those who think they discovered sliced bread for the first time a few weeks ago, we accept that there is nothing new under the sun.''

\hfill

\texttt{Thorsten on 2013-10-16 at 15:46 said:}

I'm working on understanding the Gospel in a different way, particularly the crucifixion, as something other than a glorification of pious victimhood, but it's a struggle. I understand that in the Middle Ages the story was adjusted to be more palatable to the martial character of northern barbarians and that the ``Christianity'' of the Middle Ages was less ``world-denying'' than its original form as a result. I also understand that medieval Christianity was compatible with tradition. However it is a bit odd that the more integrally compatible it became with tradition the further it moved away from its origin as a martyr myth, that it to say the more Indo-European it became. (Hellenized in terms of metaphysics, Romanized in terms of institutional organization and its enmeshing with the symbolism of empire, Germanized and Celticized in terms of chivalry and feudal ethos.)

So if we recognize this phase of ``Christianity'' as the height of its glory and accordingly agree that this is the phase that is most worth of emulation (precisely the period during which it had the least to do with victimhood and passivity) and integration into a new religion (should that ever be feasible or deemed desirable by a sufficient number of the elite) why maintain the focus on Christ (no matter how nominal)? The crucifixion seems to be not only the most inessential part of the best of what we call ``Christianity'' or the civilization of ``Christendom'', but has also been used with great efficacy by the forces of modernity to undermine the traditional values of that era.

So of course the ``religion of the future'' that I indicated as a possibility would sound strangely like Christianity (of the middle ages) but with the essential difference that it would be rectified by cutting loose the baggage of a myth that is less than optimal for the spiritual development of European man. The archetype of the triumphant hero or dragonslayer leaves no room for alternative interpretations. The myths which communicate clear cut, unambiguous mastery over oneself and the animal instincts cannot be subverted to the advantage of the Shudra.

\hfill

\texttt{Santiago on 2013-10-16 at 16:15 said:}

@Thorsten,

The point of the cross is that it's true. There is a cross to bear in our life, and we will not overcome it unless we embrace it. Thusly the Dionysian dismemberment brings forth wine, and the Odinic hanging reveals the runes.

``What you resist, persists'' as Jung put it, therefore ``resist not evil'' (Mt 5:39).

\hfill

\texttt{Michael on 2013-10-16 at 18:59 said:}

@ Thorsten. The Christian myth from an exoteric perspective is that the Divine King entered into the territory of his enemy complete disarmed, faced the full onslaught of the enemy, and overcame. The King's followers are also engaged in the cosmic battle, but the battle is spiritual, not physical.

An important component of the myth that Evola fails to note is that the King returns at the end of time in glory. All his enemies are put below his feet and he rules with a rod of iron.

If one leaves out the kingly ministry of Christ, then it is no longer Christianity.

\hfill

\end{sffamily}
\end{footnotesize}

